\chapter{绪论}

\section{研究概况}
随着国家对高等教育及科学研究投入的不断增加，高校承担的科研项目数量与经费规模呈现爆炸式增长。传统的科研经费管理模式多依赖于人工记录或功能单一的财务系统，在处理日益复杂的多源、多维数据时，往往表现出数据冗余度高、交互性差、分析维度单一等弊端。

近年来，数字化转型已成为高校治理现代化的必然选择。在此背景下，利用大数据挖掘、多维数据分析等技术构建辅助决策支持系统，已成为提升科研管理效率的关键。多维数据分析技术能够从项目进度、经费执行率、预算类别占比、历史支出趋势等多个维度对海量经费数据进行立体化切片与分析。这不仅有助于管理层实时掌握经费的全局流动情况，更能通过历史数据的积累与演化分析，为未来的经费拨付决策、预算额度核定以及科研资源优化配置提供坚实的科学依据。因此，研发一套具备多维分析能力的课题经费决策支持系统，对于实现高校科研经费的精准化、透明化管理具有重要的理论价值与实践意义。

\section{核心问题}
在当前的高校科研经费管理实践中，仍面临着诸多亟待解决的核心问题，本系统旨在针对以下痛点提供解决方案：

首先，是实时监控与风险预警能力的缺失。现有的多数系统仍属于“事后统计”范畴，管理层难以在经费执行过程中及时察觉项目逾期、余额告急或执行率严重滞后于时间进度等异常情况。本系统通过构建“智能风险标签引擎”，对项目的时间维度与经费执行维度进行实时多维比对，旨在实现从被动管理向主动预警的转变。

其次，是决策支持手段的贫乏与低效。管理人员在进行预算调剂审批或经费评估时，往往缺乏直观、多维的数据支撑，导致决策过程存在主观因素干扰。本系统引入多维数据分析模型，通过交互式可视化仪表盘，将枯燥的数字矩阵转化为趋势图、占比图及对比矩阵，显著提升了管理决策的科学性与透明度。

最后，是预算执行的全生命周期闭环管理问题。针对预算编制与实际支出的脱节现象，本系统强化了“概算-预算-支出-审计”的链条式管理。通过多维数据采集，确保了经费流向的每一笔记录均可追溯，解决了预算调剂不透明、审计追踪难的核心矛盾，从而构建起一套完整的高校项目经费决策支持体系。

\section{全文的框架结构}
本论文围绕“基于多维数据分析的高校项目经费决策支持系统”的设计与实现展开，共分为七个章节，具体安排如下：

第一章：绪论。主要阐述了课题的研究背景、研究概况以及在决策支持方面面临的核心问题，并简要介绍了全文的框架结构。

第二章：相关技术综述。对开发本系统所使用的关键技术进行了详细介绍，包括 Spring Boot 后端架构、Vue3 前端框架、MyBatis-Plus 持久层逻辑以及用于多维分析的可视化引擎 ECharts。

第三章：系统需求分析。从业务流程、功能需求以及决策支持的多维指标需求等方面进行了详细调研与建模，确定了系统必须具备的核心能力。

第四章：系统总体设计。介绍了系统的整体架构方案，包括逻辑架构、物理架构以及数据库的概念结构与逻辑结构设计，重点展示了决策分析模块的多维建模策略。

第五章：核心功能模块的实现。详细展示了后端统计服务、智能预警算法、跨组件通信机制以及多维可视化看板的编码实现细节。

第六章：系统测试与分析。通过功能测试与压力测试验证系统的稳定性，并结合实际数据对系统的决策预警准确度进行了详尽的分析与验证。

第七章：结论与展望。总结了本系统的研究工作与创新点，并针对系统目前存在的不足之处提出了未来的改进建议。
